\section{Implementation}\label{sec:impl}

\subsection{OpenMP integration}

The parallel expansion (Phase~1 of \cref{alg:par}) is implemented with a
single OpenMP pragma:

\begin{lstlisting}[language=C++, caption={The parallel expansion loop.}]
// tls[t] is thread t's private buffer (std::vector<pair<int,Dbm>>)
std::vector<std::vector<std::pair<int,Dbm>>> tls(nthreads);
#pragma omp parallel
{
    const int tid = omp_get_thread_num();
    auto& buf = tls[tid];   // private reference -- no sharing
#pragma omp for schedule(dynamic, 8)
    for (long i = 0; i < (long)frontier.size(); ++i)
        successors(W, frontier[i].first, frontier[i].second, buf, ex);
}
// Phase 2: sequential merge (single thread, no pragma needed)
std::vector<std::pair<int,Dbm>> next;
for (auto& buf : tls)
    for (auto& s : buf) { /* subsumption check + passed update */ }
frontier.swap(next);
\end{lstlisting}

\noindent\textbf{Thread-local buffers.}  The outer \texttt{tls} vector has
one entry per thread, pre-allocated once at the start of each level.  Inside
the parallel region each thread holds a reference to its own entry and never
touches another's.  The reference is resolved before the \texttt{for} loop,
so even the indirection through \texttt{tls} is private.

\noindent\textbf{Dynamic scheduling.}  The chunk size~8 in
\texttt{schedule(dynamic,8)} is chosen empirically.  The cost per state varies
widely: a state at a location with many enabled edges generates many
successors, while a state at a deadlock location generates none.  Static
assignment would create severe load imbalance.  A chunk of 8 amortises the
scheduling overhead (a few hundred nanoseconds per chunk dispatch) while
remaining fine enough for good balance.

\noindent\textbf{Graceful fallback.}  If the build has no OpenMP support
(\texttt{PTCTLMC\_HAVE\_OPENMP} not defined), the \texttt{parallel} flag is
set to \texttt{false} and the exact sequential loop of \cref{alg:seq} runs.
The two paths share the same passed-list and successor logic; only the outer
loop structure differs.

\subsection{Passed-list data structures}

The passed list for location $\ell$ is a \emph{pair} of structures:

\begin{itemize}
  \item A \texttt{std::unordered\_set<std::string>} keyed on the raw
        byte-string of the canonical DBM.  This answers the
        \emph{exact-duplicate} query (``has this precise zone been seen
        before?'') in $O(n^2)$ via hashing---the common case when two paths
        produce the same zone.
  \item A \texttt{std::vector<Dbm>} for \emph{subsumption} queries
        ($Z' \subseteq Z_p$ for some $Z_p$).  A linear scan suffices because
        LU$^+$ bounds the vector length.
\end{itemize}

In Phase~2, every candidate $(\ell', Z')$ first hits the hash set (cheap) and
only falls through to the vector scan if the exact key is absent.  This
two-tier design avoids the $O(n^2)$ inclusion test for the most frequent case.

\subsection{Zone operations}

All DBM operations (intersection, reset, delay, LU$^+$ extrapolation,
canonical closure, inclusion) are provided by the \textsc{HDBM} library
bundled with \tool{}.  Each DBM is stored as a flat
\texttt{std::vector<raw\_t>} where \texttt{raw\_t} is a 32-bit integer
encoding both the bound value and its strictness in the MSB.  Canonical form
is maintained after every mutating operation.
